% =============================================================================
% Appendix A — Solutions
% =============================================================================
\section{Solutions to Selected Exercises}\label{app:solutions}

Solutions are given for the exercises whose answers are not immediate. Where an
exercise asks for a design or an argument, the solution indicates what a complete
answer must contain rather than supplying prose to copy.

\subsection*{Section 1}

\paragraph{1.1}
(a) Arithmetic --- falsified only by a computational error.
(b) Modelling --- falsified by evidence that stages interact, e.g.\ an error rate
that is not the product of per-stage factors.
(c) Measurement --- falsified by an independent determination in the same
background and reporter.
(d) Modelling --- falsified by a tumour with high polymerase error rate and low
burden, or the reverse; Section~\ref{sec:mmr} argues exactly this.

\paragraph{1.2}
A reply should make two points: that the gate establishes the conclusion follows
from the hypotheses and nothing about whether the hypotheses hold of a cell; and
that this is not a criticism of formalisation, whose value is that a hypothesis
precise enough to formalise is precise enough to refute.

\subsection*{Section 2}

\paragraph{2.1}
By Proposition~\ref{prop:gain} with $k = 1$, the ratio
$\gamma^{m+1}/\gamma^{m} = \gamma$ for every $m$. Returns are constant in the
logarithm because the stages are statistically independent and the errors
multiply; most engineered schemes share information between stages, so later
stages have less to contribute.

\paragraph{2.2}
With $\mathcal{E}_2 = \gamma^2 + c$ and $c > 0$ we can have
$\mathcal{E}_2 > \mathcal{E}_1 = \gamma$ whenever $c > \gamma - \gamma^2$. So a
two-stage enzyme is worse than a one-stage one: monotonicity fails. The
independence assumption is therefore not cosmetic --- it is what makes
``more proofreading is never worse'' true.

\paragraph{2.3}
$\gamma = e^{-\Delta\Delta G/k_BT} = 10^{-6}$ requires
$\Delta\Delta G = 6\ln 10 \approx 13.8\,k_BT$. At $310$\,K,
$k_BT \approx 2.58\ \mathrm{kJ\,mol^{-1}}$, so
$\Delta\Delta G \approx 35.6\ \mathrm{kJ\,mol^{-1}}$. A single Watson--Crick
hydrogen bond contributes of order $5$--$10\ \mathrm{kJ\,mol^{-1}}$, and the
difference between a correct and incorrect pair is smaller still, so no plausible
base-pairing chemistry supplies it. This is precisely why equilibrium
discrimination is insufficient and kinetic proofreading is needed.

\subsection*{Section 3}

\paragraph{3.1}
Apply Theorem~\ref{thm:psr1} with the ambient set $\kept$ and the subset
$\emptyset$: it gives $\errate(\kept) \le \errate(\emptyset) = \inserr$, since the
empty product is $1$. Hence PSR$_2$.

\paragraph{3.2}
(a) Multiplicativity fails: the error rate is no longer a product over
independent stages. The ceiling need not hold, and this is the one escape route
the no-go theorem leaves open.
(b) No hypothesis of Assumption~\ref{model:class} fails, but the mutant is not a
\emph{stage-removal} mutant --- $\inserr$ has changed. This is the case PSR$_3$
describes, and it is the rescue the Xing measurement forbids.
(c) The range hypothesis $\escape_i \in [0,1]$ fails, and the ceiling genuinely
fails with it: the Lean file's negative control exhibits this.

\paragraph{3.3}
A complete answer defines refutation of a family as the non-existence of any
$\theta \in \Theta$ consistent with the observation, and refutation of a member
as inconsistency for one particular $\theta$. Enlarging $\Theta$ weakens content
because the set of observations inconsistent with \emph{every} member shrinks;
in the limit of a sufficiently rich family, nothing is excluded and the family
makes no prediction.

\subsection*{Section 4}

\paragraph{4.1}
PSR$_3$ with $\rho = 2300/44 \approx 52.3$ gives
$\inserr' \ge 52.3\,\inserr$: the insertion/extension error would have to exceed
wild type's by a factor of about $52$. Xing's purified measurement forbids it,
because with no repair present \texttt{P301R} makes \emph{fewer} errors than
exonuclease-null, so its intrinsic error is not elevated at all.

\paragraph{4.2}
It establishes only that this allele is consistent with the model class --- it
does not confirm the model, because every allele that is \emph{not} a mutator is
equally consistent. A test that cannot fail for non-mutators provides no evidence
when it passes. (Compare Exercise~4.4.)

\paragraph{4.4}
If the statistic is forced to its extreme value by a symmetry of the pipeline,
then every shuffle produces a value at least as extreme, so the $p$-value is
bounded by $1/(N+1)$ regardless of the data. It measures the pipeline, not the
hypothesis. The rule violated: \emph{a control must be able to come out
non-significant}. Before reporting any test, construct the observation that would
fail it; if none exists, the test carries no evidence.

\subsection*{Section 5}

\paragraph{5.1}
$\mathrm{esc}(3, 0) = 3 \le 7 = \mathrm{esc}(7, 0)$ (neither saturated);
$\mathrm{esc}(3, 5) = 0 \le 2 = \mathrm{esc}(7, 5)$ (the mutant is fully repaired);
$\mathrm{esc}(3, 9) = 0 \le 0 = \mathrm{esc}(7, 9)$ (both fully repaired).
The ordering holds in every case, which is Theorem~\ref{thm:norev}.

\paragraph{5.2}
$\partial_p\left[p^2/(p+c)\right] = p(p + 2c)/(p+c)^2 \ge 0$ and
$\partial_c\left[p^2/(p+c)\right] = -p^2/(p+c)^2 \le 0$. Both monotonicities
hold, so Theorems~\ref{thm:norev} and~\ref{thm:forced} go through unchanged. The
conclusion depends only on the \emph{signs} of these derivatives, not on the
hard threshold, which is why the caricature is acceptable.

\paragraph{5.4}
In a null background $c = 0$ for both enzymes, so
$\mathrm{esc}(p, 0) = p$ and the comparison reduces to intrinsic production ---
which is already known from the purified assay and cannot discriminate the
hypotheses. A hypomorphic background gives $0 < c < c_{\mathrm{null}}$, where
titration (capacity consumed in proportion to load) and interference (capacity
reduced independently of load) predict different dependences of escape on $p$.

\subsection*{Section 6}

\paragraph{6.1}
For $g = 4$, $\mu = 10^{-2}$: the sum $\sum_{j=0}^{3} \mu 2^{j} 2^{3-j} =
4\mu\cdot 2^{3} = 0.32$, and $\tfrac{1}{2}\mu g N = \tfrac{1}{2}(10^{-2})(4)(16)
= 0.32$. They agree.

\paragraph{6.4}
Both cultures have the same $g$ and $N$; by \eqref{eq:ldmean} the means are
$\tfrac12\mu_1 gN$ and $\tfrac12\mu_2 gN$, whose ratio is $\mu_2/\mu_1 = \rho$.
Every factor depending on the growth protocol cancels, which is why the ratio is
the robust statistic and the individual rates are not.

\subsection*{Section 7}

\paragraph{7.2}
Two errors. First, the categories overlap, so pathogenic-over-total is not a
proportion of anything: the denominator includes records also counted in the
numerator's complement, and the categories sum past the total. Second, even a
correctly computed proportion of \emph{submitted} records is not a probability
for a \emph{new} variant, because ClinVar is ascertainment-biased toward variants
someone had a reason to submit.

\subsection*{Section 8}

\paragraph{8.1}
Unrolling \eqref{eq:cgr}, $x_i = \sum_{k\ge0} 2^{-(k+1)} v_{b_{i-k}}$. The
coefficient of $v_{b_{i-k}}$ is $2^{-(k+1)}$. If two positions share their last
$k$ bases, the difference of their embeddings is supported on terms of index
$> k$, hence bounded by $\sum_{j>k} 2^{-(j+1)}\mathrm{diam} = 2^{-(k+1)}\mathrm{diam}$.

\paragraph{8.2}
At the filtration limit the complex is connected for any sufficiently dense
sample of a connected attractor, so $\beta_0 = 1$ carries no information about
the sequence. It becomes informative only at $\epsilon$ below the typical
nearest-neighbour spacing, where $\beta_0$ counts clusters --- and then it is
measuring sampling density, not biology.

\subsection*{Section 9}

\paragraph{9.1}
Theorem~\ref{thm:trade} with $\Delta\Tw = -1$ gives $\Delta\Wr = +1$. Physically
the compensating writhe appears as a crossing of the axis --- a supercoil --- in
the flanking DNA.

\paragraph{9.2}
$\Lk$ counts signed crossings of two closed curves and is an integer for
topological reasons. $\Tw$ and $\Wr$ are integrals and individually real; their
sum is constrained to be the integer $\Lk$ by Theorem~\ref{thm:cwf}. If the
ribbon is not closed, $\Lk$ is undefined and the decomposition fails --- which is
why a nicked plasmid has no superhelical density.

\paragraph{9.3}
$95$ turns per $1000$ bp at $50$ bp/s is $95 \times (50/1000) = 4.75$ turns per
second.

\subsection*{Section 10}

\paragraph{10.1}
$\frac{1}{5}\cdot\frac{62500}{9} = \frac{62500}{45} = \frac{12500}{9}$. Verifying
it confirms a property of the model, and the model has been refuted; a correct
consequence of a false premise is not evidence for anything.

\paragraph{10.3}
Claims that would change: anything deriving a burden or rate from stage counts.
Claims that would not: the live database records of Section~\ref{sec:variants},
the topology of Section~\ref{sec:dnatop}, the clinical decision rule. The
distinguishing feature is whether the claim takes the mechanism as an input or
measures its consequence directly.
